\title{Geometric deep learning on molecular representations}
\begin{document}
\maketitle

\begin{abstract}
Review Article https://doi.org/10.1038/s42256-021-00418-8 1Department of Chemistry and Applied Biosciences, RETHINK, ETH Zurich, Zurich, Switzerland. 2Department of Biomedical Engineering, Eindhoven University of Technology, Eindhoven, the Netherlands. 3ETH Singapore SEC Ltd, Singapore, Singapore. 4These authors contributed equally: Kenneth Atz, Francesca Grisoni. ✉e-mail: f.grisoni@tue.nl; gisbert@ethz.ch R ecent advances in deep learning, an instance of neural network-based artificial intelligence (AI)1,2, have led to pio- neering applications in the molecular sciences, for example, in drug discovery3, quantum chemistry4 and structural biology5. Two characteristics of deep learning render it promising when applied to molecules. First, deep learning methods can cope with ‘unstruc- tured’ data representations such as text sequences6, speech signals7, images8 and graphs9. This ability seems particularly useful for molecular systems, for which chemists have developed molecular representations that capture molecular properties at varying levels of abstraction (Fig. 1). Second, deep learning can perform feature extraction (or feature learning) from the input data; that is, produce data-driven features from the input representation. These two char- acteristics of deep learning complement ‘classical’ machine learning applications (such as quantitative structure–activity relationships (QSAR)), in which molecular features (that is, ‘molecular descrip- tors’10) are encoded a priori wi
\end{abstract}

\section{Recovered Text}
Review Article
https://doi.org/10.1038/s42256-021-00418-8
1Department of Chemistry and Applied Biosciences, RETHINK, ETH Zurich, Zurich, Switzerland. 2Department of Biomedical Engineering, Eindhoven
University of Technology, Eindhoven, the Netherlands. 3ETH Singapore SEC Ltd, Singapore, Singapore. 4These authors contributed equally: Kenneth Atz,
Francesca Grisoni. ✉e-mail: f.grisoni@tue.nl; gisbert@ethz.ch
R
ecent advances in deep learning, an instance of neural
network-based artificial intelligence (AI)1,2, have led to pio-
neering applications in the molecular sciences, for example, in
drug discovery3, quantum chemistry4 and structural biology5. Two
characteristics of deep learning render it promising when applied
to molecules. First, deep learning methods can cope with ‘unstruc-
tured’ data representations such as text sequences6, speech signals7,
images8 and graphs9. This ability seems particularly useful for
molecular systems, for which chemists have developed molecular
representations that capture molecular properties at varying levels
of abstraction (Fig. 1). Second, deep learning can perform feature
extraction (or feature learning) from the input data; that is, produce
data-driven features from the input representation. These two char-
acteristics of deep learning complement ‘classical’ machine learning
applications (such as quantitative structure–activity relationships
(QSAR)), in which molecular features (that is, ‘molecular descrip-
tors’10) are encoded a priori with rule-based algorithms. This capa-
bility of multi-layered neural networks to learn from unstructured
data and extract higher-order molecular features has led to numer-
ous applications of deep learning in the molecular sciences.
Geometric deep learning (GDL) is an emerging concept of
AI. GDL is an umbrella term encompassing emerging techniques
that generalize neural networks to Euclidean and non-Euclidean
domains, such as graphs, manifolds, meshes or string representa-
tions9. In general, GDL includes approaches that incorporate a
geometric prior—that is, information on the structure space and
symmetry properties of input signals, such as representations of
molecular structures. The geometric prior is leveraged to improve
the quality of the model, for example its predictive accuracy.
Although GDL has been increasingly applied to molecular model-
ling4,11–13, its full potential in the field is still untapped.
The aim of this Review is to (1) provide a structured and har-
monized overview of prominent applications of GDL in molecu-
lar systems, (2) delineate the main research directions in the field
and (3) attempt a critical forecast of the potential future impact of
GDL. Three fields of application are highlighted, namely drug dis-
covery, quantum chemistry and computer-aided synthesis planning
(CASP). A glossary of selected terms can be found in Box 1.
Principles of GDL
The term GDL was coined in 20179. Although GDL was originally
used for methods applied to non-Euclidean data9, it now extends
to all deep learning methods that incorporate geometric priors14.
Symmetry is a crucial concept in GDL, as it encompasses the prop-
erties of the system with respect to manipulations (transformations),
such as the ones covered by the Euclidean group E(3) (Box 2). Other
relevant transformations for molecules include scale separation (for
example, coarse and fine graining of a grid) or permutation (that is,
different ordering of the atoms in a molecule).
Symmetry is often recast in terms of invariance and equivariance
to express the behaviour of any mathematical function with respect
to a transformation T  (for example rotation, translation, reflection
or permutation) of an acting symmetry group15. Here, the mathe-
matical function is a neural network F applied to a given molecular
input X. F(X) can therein transform equivariantly, invariantly or
non-equivariantly with respect to T , as described below:
•
Equivariance. A neural network F applied to an input X
is equivariant to a transformation T  if the transforma-
tion of X commutes with the transformation of F(X) via a
transformation T ′ of the same symmetry group, such that:
F(T (X)) = T ′F(X). Neural networks are therefore equivari-
ant to the actions of a symmetry group on their inputs if, and
only if, each layer of the network ‘equivalently’ transforms under
any transformation of that group.
•
Invariance. Invariance is a special case of equivariance where
F(X) is invariant to T  if T ′ is the trivial group action (that is,
identity): F(T (X)) = T ′F(X) = F(X).
•
Non-equivariance. F(X) is neither equivariant nor invariant to
T  when the transformation of X does not commute with the
transformation of F(X): F(T (X)) ̸= T ′F(X).
The symmetry properties of a neural network architecture vary
depending on the network type and the symmetry group of inter-
est, and are individually discussed in the following sections. Readers
can find in-depth treatments of equivariance and group equivariant
layers in neural networks elsewhere16,17.
Geometric deep learning on molecular
representations
Kenneth Atz   1,4, Francesca Grisoni   1,2,4 ✉ and Gisbert Schneider   1,3 ✉
Geometric deep learning (GDL) is based on neural network architectures that incorporate and process symmetry information.
GDL bears promise for molecular modelling applications that rely on molecular representations with different symmetry prop-
erties and levels of abstraction. This Review provides a structured and harmonized overview of molecular GDL, highlighting
its applications in drug discovery, chemical synthesis prediction and quantum chemistry. It contains an introduction to the
principles of GDL, as well as relevant molecular representations, such as molecular graphs, grids, surfaces and strings, and
their respective properties. The current challenges for GDL in the molecular sciences are discussed, and a forecast of future
opportunities is attempted.
Nature Machine Intelligence | VOL 3 | December 2021 | 1023–1032 | www.nature.com/natmachintell
1023
Review Article
NaTuRe MachIne In TeLLIGence
The concept of equivariance and invariance can also be used in
reference to the molecular features obtained from a given molecular
representation (X), depending on their behaviour when a transfor-
mation is applied to X. For instance, many molecular descriptors are
invariant to the rotation and translation of the molecular represen-
tation by design10, for example, the Moriguchi octanol–water parti-
tioning coefficient18, which relies only on the occurrence of specific
molecular substructures for calculation. The symmetry properties
of molecular features extracted by a neural network depend on the
symmetry properties of both the input molecular representation
and the utilized neural network.
Many relevant molecular properties (for example, equilibrium
energies, atomic charges or physicochemical properties such as per-
meability, lipophilicity or solubility) are invariant to certain sym-
metry operations (Box 2). For applications in chemistry it can thus
be desirable to design neural networks that transform equivariantly
under the actions of predefined symmetry groups. Exceptions occur
if the targeted property changes following a symmetry transforma-
tion of the molecules (for example, chiral properties that change
under inversion of the molecule or vector properties that change
under rotation of the molecule). In such cases, the inductive bias
(learning bias) of equivariant neural networks would not allow for
the differentiation of symmetry-transformed molecules.
Whereas neural networks can be considered as universal func-
tion approximators19, incorporating prior knowledge such as rea-
sonable geometric information (geometric priors) has evolved as a
core design principle of neural network modelling14. By incorpo-
rating geometric priors, GDL allows the quality of the model to be
increased and bypasses several bottlenecks related to the need to
force the data into Euclidean geometries (for example, by feature
engineering). Moreover, GDL provides novel modelling opportuni-
ties, such as data augmentation in low-data regimes.
Molecular GDL
The application of GDL to molecular systems is challenging, partly
because there are multiple valid ways of representing the same
molecular entity. Molecular representations can be categorized on
the basis of their different levels of abstraction and the physicochem-
ical and geometrical aspects they capture. Importantly, all of these
representations are models of the same reality and are thus suitable
for some purposes, not for others20. GDL provides the opportunity
to experiment with different representations of the same molecule
and leverages their intrinsic geometrical features to improve the
quality of the model. Moreover, GDL has repeatedly proved useful in
providing insights into relevant molecular properties for the task at
hand thanks to its feature extraction (feature learning) capabilities.
In the following sections, we delineate the most prevalent molecu-
lar GDL approaches (see Table 1 for an overview) and their applica-
tions in chemistry, grouped according to the respective molecular
representations used for deep learning: molecular graphs and point
clouds, grids, surfaces and string notations. To avoid confusion with
other usages of the word ‘representation’ in deep learning, we will use
the term ‘feature’ whenever we refer to any numerical description of
molecules, obtained either with rule-based algorithms (molecular
descriptors) or learned (extracted) by neural networks.
Learning on molecular graphs. Molecular graphs. Graphs are
among the most intuitive ways to represent molecular struc-
tures21. Any molecule can be thought of as a mathematical graph
G = (V, E) (where G refers to the molecular graph, V to the set of
vertices and E to the set of adjacent edges), whose vertices (vi ∈V)
represent atoms, and whose edges (eij ∈E) constitute their connec-
tion (Fig. 3a). In many deep learning applications, molecular graphs
can be further characterized by a set of vertex and edge features.
3D molecular structures are commonly represented as 3D graphs
G3D = (V, E, R) or 3D point clouds P3D = (V, R), where in addi-
tion to vertex and potential edge features (vi ∈V and eij ∈E,
respectively), information on the vertex position in a 3D coordinate
system (ri ∈R) is encoded (where R refers to the set of coordi-
nates). 3D point clouds are a special case of 3D graphs in which no
edges or edge features are predefined, and the interacting vertices
are either specified through a given radius or the complete edge set.
GNNs. Deep learning methods that handle graphs as input are com-
monly referred to as GNNs21–23. When applied to molecular graphs,
GNNs perform feature extraction by progressively aggregating
information from atoms and their molecular environments (Fig. 2a;
refs. 24,25). Different architectures of GNNs have been introduced14,26,
the most popular of which fall under the umbrella term of message
passing neural networks4,25. Such networks iteratively update the
vertex features of the lth network layer (vl
i →vl+1
i
) via message pass-
ing operations, using at least two learnable functions ψ and ϕ, and a
local permutation-invariant aggregation operator ⨁ (for example,
sum, mean or maximum): vl+1
i
= ϕ
(
vl
i, ⊕
j∈N (i)ψ (vl
i, vl
j
))
.
Since their introduction as a means to predict quantum chemi-
cal properties of small molecules at the DFT level4, GNNs have
found many applications in quantum chemistry27,28, drug discov-
ery29–31, CASP32,33 and molecular property prediction34,35. When
applied to quantum chemistry tasks, GNNs can use 3D informa-
tion by including radial and angular information in the edge fea-
tures of the graph28,36,37, which has proved beneficial for estimating
quantum chemical energies and other properties of equilibrium and
non-equilibrium molecular conformations, as in the case of SchNet38
and PaiNN37. SchNet-like architectures were used to predict quan-
tum mechanical wavefunctions in the form of Hartree–Fock and
DFT density matrices39, and differences in quantum properties
obtained by DFT and coupled cluster level-of-theory calculations40.
GNNs for molecular property prediction have been shown to
outperform models based on human-engineered molecular descrip-
tors for certain biologically relevant properties41. Although includ-
ing 3D information into molecular graphs generally improved the
prediction of drug-relevant properties, there was no marked differ-
ence between using single or multiple molecular conformations for
network training42. GNN models seem suitable for explainable AI43,
CC1(C)[C@H](C(O)=O)N2[C@@H](CC2)S1
c
a
b
d
e
Fig. 1 | Exemplary molecular representations for a selected molecule.
The penam substructure of penicillin is shown. a, 2D depiction (Kekulé
structure). b, Molecular graph (2D) composed of vertices (atoms) and edges
(bonds). c, Simplified molecular input line entry system (SMILES) string84, in
which atom type, bond type and connectivity are specified by alphanumeric
characters. d, 3D graph composed of vertices (atoms), their positions (x,
y, z coordinates) in 3D space and edges (bonds). e, Molecular surface
represented as a mesh coloured according to the respective atom types.
Nature Machine Intelligence | VOL 3 | December 2021 | 1023–1032 | www.nature.com/natmachintell
1024
Review Article
NaTuRe MachIne In TeLLIGence
which was used to interpret models predicting molecular properties
of preclinical relevance44 and quantum chemical properties45. GNNs
have also been used for molecular fingerprinting, albeit with mixed
results22,46,47.
It is also possible to generate molecules de novo with GNNs48–51,
for example by adding graph vertices (atoms) and edges (bonds) to
an initial vertex48, or in combination with variational autoencod-
ers49,50,52 and reinforcement learning51,53,54. Finally, GNNs have been
applied to CASP32,33,55. However, current approaches are limited to
chemical transformations in which one bond is removed between
the products and the reactants.
Graph isomorphism is an active research topic in graph-based
machine learning56–59. When considering only topological infor-
mation (that is, the number of edges and vertices, and vertex adja-
cency), message passing neural networks have been shown to be
unable to distinguish selected examples of non-isomorphic graphs25
(that is, they are learning identical features for two different graphs).
This aspect is relevant for molecular pattern recognition (for exam-
ple property prediction) when two different molecules should result
in distinctive learned features. Encoding angular and radial features
allows GNNs to capture geometric shapes (such as pentagons or
hexagons), thereby enabling them to distinguish between a broader
spectrum of non-isomorphic molecular graphs28.
One recent area of development of graph-based methods is SE(3)-
and E(3)-equivariant GNNs (equivariant message passing networks),
which deal with the absolute coordinate systems of 3D graphs60
(Fig. 2b) by exploiting relevant Euclidean symmetries (Box 2). This
ability is enabled by defining a respective network architecture. By
using E(3)-equivariant61 and SE(3)-equivariant60,62 convolutions, such
networks have demonstrated high accuracy for predicting several
quantum chemical properties, including energies37,62–65, interatomic
potentials for molecular dynamics simulations61,64,66 and wavefunc-
tions67. Such SE(3)-equivariant neural networks do not commute
with reflections of the input (that is, they are non-equivariant to
reflections), and thereby enable SE(3)-equivariant models to dis-
tinguish between stereoisomers of chiral molecules including enan-
tiomers60. By contrast, E(3)-equivariant neural networks transform
equivariantly with reflections, which allows E(3)-equivariant mod-
els to distinguish between diastereomers but not enantiomers. SE(3)
neural networks are computationally expensive due to their use of
spherical harmonics and Wigner D-matrices to compute learnable
weight kernels60. E(3)-equivariant neural networks are computation-
ally more efficient and have been shown to perform as well as, or bet-
ter than, SE(3)-equivariant networks; for example, for the modelling
of quantum chemical properties and dynamic systems37,61. Recent
applications of SE(3)-equivariant neural networks led to scoring
Box 1 | |Glossary of selected terms
CoMFA and CoMS IA. Comparative molecular field analysis
(CoMFA)72 and comparative molecular similarity indices analysis
(CoMSIA)73 are 3D QSAR methods developed in the 1980s and
1990s in which 3D grids are used to capture the distributions of
molecular features (such as the steric, hydrophobic and electro-
static properties). The molecular descriptors obtained serve as in-
puts to a regression model for quantitative bioactivity prediction.
Convolution. An operation within a neural network that
transforms a feature space into a new feature space and thereby
captures the local information found in the data. Convolutions
were originally introduced for pixels in images145, but the term
is also used for neural network architectures covering a variety
of data structures such as graphs, point clouds, spheres, grids or
manifolds.
Density functional theory (DFT). A quantum mechanical
modelling approach used to investigate the electronic structure of
molecules.
Data augmentation. Artificial increase of the data volume available
for model training, often achieved by leveraging symmetrical
properties of the input data that are not captured by the model (for
example, rotation or permutation).
Feature. An individually measurable or computationally obtainable
characteristic of a given sample (for example, a molecule), often in
the form of a scalar. In this Review, the term refers to a numeric
value characterizing a molecule. Such molecular features can be
computed with rule-based algorithms (molecular descriptors)
or generated automatically by deep learning from a molecular
representation (‘hidden’ or ‘learned’ features).
Geometric prior. An inductive bias incorporating information
on the symmetric nature of the system of interest into the neural
network architecture. Also known as a symmetry prior.
Graph isomorphism. Two graphs are considered isomorphic if
there exists a one-to-one correspondence between their vertices
(bijection) that preserves the corresponding adjacency and
non-adjacency. In other words, isomorphic graphs have the same
number of elements (vertices and edges) and the same vertex
adjacency.
Inductive bias. A set of assumptions that a learning algorithm
(such as a neural network) uses to learn the target function and to
make predictions on previously unseen data points.
One-hot encoding. Method for representing categorical variables
as numerical arrays by obtaining a binary variable (0, 1) for each
category. It is often used to convert sequences (for example,
SMILES strings) into numerical matrices that are suitable inputs
to and/or outputs of deep learning models (such as chemical
language models).
Quantitative Structure-Activity Relationship QSAR. Machine
learning techniques that aim to find an empirical relationship
between the molecular structure (usually encoded as molecular
descriptors) and experimentally determined molecular properties,
such as physicochemical properties or bioactivity.
Reinforcement learning. A technique used to steer the output
of a machine learning algorithm towards user-defined regions of
optimality via a predefined reward function146.
Transfer learning. The transfer of knowledge from an existing
deep learning model to a related task for which fewer training
samples are available147.
Unstructured data. Data that are not arranged as vectors of
features. Examples of unstructured data include graphs, images
and meshes. Molecular representations are typically unstructured,
whereas numerical molecular descriptors (such as molecular
properties and molecular ‘fingerprints’) are examples of structured
data.
Voxel. Element of a regularly spaced three-dimensional (3D) grid
(equivalent to a pixel in two-dimensional (2D) space).
Nature Machine Intelligence | VOL 3 | December 2021 | 1023–1032 | www.nature.com/natmachintell
1025
Review Article
NaTuRe MachIne In TeLLIGence
functions for folded 3D structures of ribonucleic acid11. Equivariant
message passing networks have also been applied to predicting the
quantum mechanical wavefunctions of nucleus- and electron-based
representations68–70. However, such networks are limited at present
to small molecular systems because of the large size of the learned
matrices, which scale quadratically with the number of electrons in
the system.
Learning on grids. Grids capture the properties of a system at regu-
larly spaced intervals. Depending on the number of dimensions
included in the system, grids can be 1D (for example, sequences), 2D
(RGB images), 3D (cubic lattices) or higher-dimensional. Grids are
defined by a Euclidean geometry and can be considered as a graph
with a special adjacency, where (1) each vertex has an identical num-
ber of adjacent edges (that is, identical neighbourhood structure),
and is therefore indistinguishable from all other vertices structure
wise, and (2) the vertices have a fixed ordering that is defined by the
spatial dimensions of the grid14. These two properties render local
convolutions applied to a grid inherently permutation invariant, and
provide a strong geometric prior for translation equivariance (for
example, by weight sharing in convolutions) as well as scale separa-
tion (for example, by pooling operations)14. These grid properties
have critically determined the success of CNNs in computer vision8.
Molecular grids. Molecules can be represented as grids in different
ways. 2D grids (for example, molecular structure drawings, Kekulé
structures) are generally more useful for visualization than predic-
tion, with few exceptions71. In analogy to some popular pre-deep
learning approaches, such as CoMFA72 and CoMSIA73, 3D grids
are often used to capture the spatial distribution of the properties
within one (or more) molecular conformer. Such representations
are then used as input to the 3D CNNs. 3D CNNs are character-
ized by a greater resource efficiency than equivariant GNNs, which
until now have mainly been applied to molecules with fewer than
approximately 1,000 atoms. Thus, 3D CNNs have often been the
method of choice when the protein structure has to be considered—
for example, for protein–ligand binding affinity prediction74,75 or
active site recognition76.
Learning on molecular surfaces. Molecular surfaces can be defined
by the surface enclosing the 3D structure of a molecule at a certain
distance from each atom centre. Each point on such a continuous
Box 2 | | Euclidean symmetries in molecular systems
Molecular systems (and 3D representations thereof) can be con-
sidered as objects in Euclidean space. In such a space, one can
apply several symmetry operations (transformations) that are
performed with respect to three symmetry elements (that is,
line, plane, point) and rigid (that is, they preserve the Euclidean
distance between all pairs of atoms (isometry)). The Euclidean
transformations are as follows:
•
Rotation. Movement of an object with respect to the radial
orientation to a given point.
•
Translation. Movement of every point of an object by the
same distance in a given direction.
•
Reflection. Mapping of an object to itself through a point
(inversion), a line or a plane (mirroring).
All three transformations and their arbitrary finite
combinations are included in the Euclidean group [E(3)]. The
special Euclidean group [SE(3)] comprises only translations and
rotations.
Molecules are always symmetric in the SE(3) group; that is, their
intrinsic properties (for example, biological and physicochemical
properties and equilibrium energy) are invariant to coordinate
rotation and translation and combinations thereof. Several
molecules are chiral: some of their (chiral) properties depend on
the absolute configuration of their stereogenic centres, and are
thus non-invariant to molecule reflection. Chirality plays a key
role in chemical biology; relevant examples of chiral molecules
are DNA, and several drugs whose enantiomers exhibit markedly
different pharmacological and toxicological properties148.
Table 1 | Summary of selected GDL approaches for molecular
modelling
GDL approach
Molecular
representation(s)
Applications
Graph neural
network (GNN)
2D and 3D molecular
graphs and 3D point
clouds
Molecular property
prediction in drug
discovery29,44 and in
quantum chemistry for
energies61–63, forces37,61,64 and
wavefunctions 67, CASP33,144
and generative molecular
design 51,52.
3D convolutional
neural network
(3D CNN)
3D grid
Structure-based drug design
and property prediction74,75.
Mesh
convolutional
neural network
(geodesic CNN or
3D GNN)
Surfaces (meshes)
encoded as a 2D grids
or 3D graphs
Protein–protein interaction
prediction and ligand-pocket
fingerprinting12.
Recurrent neural
network (RNN)
String notations (1D
grid)
Generative molecular
design13,108, synthesis
planning87, protein structure
prediction112 and prediction
of properties in drug
discovery110,111.
Transformers
String notations
encoded as a graphs
Synthesis planning117,
prediction of reaction
yields121, generative molecular
design125, prediction of
properties in drug discovery123
and protein structure
prediction5,126.
For each approach, the utilized molecular representation(s) and selected applications are reported.
Original
Rotation
Translation
Reflection (mirroring)
Reflection (inversion)
Nature Machine Intelligence | VOL 3 | December 2021 | 1023–1032 | www.nature.com/natmachintell
1026
Review Article
NaTuRe MachIne In TeLLIGence
surface can be further characterized by its chemical (for example,
hydrophobic, electrostatic) and geometric (such as local shape, cur-
vature) features. From a geometric perspective, molecular surfaces
are considered as 3D meshes; that is, a set of polygons (faces) that
describe how the mesh coordinates exist in the 3D space77. Their ver-
tices can be represented by a 2D grid structure (where four vertices
on the mesh define a pixel) or by a 3D graph structure. The grid- and
graph-based structures of meshes enable applications of 2D CNNs,
geodesic CNNs and GNNs to learn on mesh-based molecular sur-
faces. Recently, geodesic (2D) CNNs have been applied to learn on
mesh-based representations of protein surfaces to predict protein–
protein interactions and recognize corresponding binding sites12.
This approach generated data-driven fingerprints that are relevant
for specific biomolecular interactions. Approaches like 2D CNNs
applied to meshes come with certain limitations, such as the need for
rotational data augmentation (due to their non-equivariance to rota-
tions) and the need to enforce a homogeneous mesh resolution (that
is, uniform spacing of all the points in the mesh). Recently introduced
GNNs for mesh-based representations have been shown to incorpo-
rate rotational equivariance into their network architecture and allow
for heterogeneous mesh resolution78. Such GNNs are computation-
ally efficient and have potential for modelling macromolecular struc-
tures; however, they have not yet found applications in molecular
systems. Other studies have used 3D voxel-based surface representa-
tions of (macro)molecules as inputs to 3D CNNs, for example, for
protein–ligand affinity79 and protein binding-site80 prediction.
Learning on string representations. Molecular strings. Molecules
can be represented as molecular strings—that is, linear sequences
of alphanumeric symbols. Molecular strings were originally devel-
oped as manual ciphering tools to complement systematic chemi-
cal nomenclature81 and later became suitable for data storage and
retrieval. Some of the most popular string-based representations are
the Wiswesser line notation82, the International Chemical Identifier
(InChI)83 and SMILES84.
Each type of linear representation can be considered as a ‘chemi-
cal language’. In fact, such notations possess a defined syntax; that
is, not all possible combinations of alphanumeric characters will
lead to a ‘chemically valid’ molecule. Furthermore, these notations
possess semantic properties: depending on how the elements of the
string are combined, the corresponding chemical compound will
have different physicochemical and biological properties. These
characteristics extend the deep learning methods developed for lan-
guage and sequence modelling to the analysis of molecular strings
for ‘chemical language modelling’85,86. Before the actual model
building, molecular strings are usually preprocessed as either 1D
molecular grids or as molecular graphs, thereby inheriting the geo-
metric prior of the corresponding representation.
SMILES strings—in which letters are used to represent atoms,
whereas symbols and numbers encode bond types, connectiv-
ity, branching and stereochemistry (Fig. 3a)—are frequently used
for sequence-based deep learning13,87. Different SMILES strings
can be obtained for a given molecule depending on the starting
non-hydrogen atom and direction in the molecular graph (Fig. 3a),
while preserving the chemical information captured. Several algo-
rithms have been proposed to obtain ‘canonical’ SMILES88,89; that is,
strings that are invariant to the order in which atoms are presented.
Whereas several other string representations have been tested
in combination with deep learning, such as InChI90, DeepSMILES91
and self-referencing embedded strings (SELFIES)92, SMILES strings
remain the de facto representation of choice for chemical language
modelling, probably due to their widespread usage in chemistry
and certain advantages for generative modelling93. String repre-
sentations such as SELFIES are receiving increasing attention from
the scientific community and are in continuous development. The
next section introduces the most prominent chemical language
Feature labelling
= (ν,ε,   )
r  = (ρ ,θ ,φ )
e
v  ∈       υ
Feature updates
Aggregation
Molecular property
Atomic property
Bond property
Message passing
a
b
Molecular property
Atomic property
Equivariant message passing
Feature labelling
Feature updates
Aggregation
vT
vT
e
y
y
y
x
y
x
y
v  ∈
∈
= (ν,ε)
,
,
,
υ
Fig. 2 | Deep learning on molecular graphs. a, Message passing GNNs applied to 2D molecular graphs. A 2D molecular graph G = (V, E) is shown with its
labelled vertex (atom) features (vi ∈Rdv) and edge (bond) features (eij ∈Rde), where R refers to the set of real numbers, and dv and de (as a superscript
of R) refer to the d-dimensional set of real numbers for vertex and edge features, respectively. Vertex features are updated by iterative message passing
for a defined number of time steps T across each pair of vertices vi and vj, connected via an edge ej,i. After the last message passing convolution, the
final vertex vt
i can be (1) mapped to a bond (yij) or atom (yi) property or (2) aggregated to form molecular features (that can be mapped to a molecular
property y). b, E(3)-equivariant message passing GNNs applied to 3D molecular graphs. 3D graphs G3D = (V, E, R) that are labelled with atom features
(vi ∈Rdv), their absolute coordinates in 3D space (ri ∈R3, where ri refers to the spherical coordinates ρ, θ and ϕ) and their edge features (eij ∈Rde) are
shown. Iterative spherical convolutions are used to obtain data-driven atomic features (vt
i), which can be mapped to atomic properties or aggregated and
mapped to molecular properties (yi and y, respectively).
Nature Machine Intelligence | VOL 3 | December 2021 | 1023–1032 | www.nature.com/natmachintell
1027
Review Article
NaTuRe MachIne In TeLLIGence
modelling methods, along with selected examples of their applica-
tion to chemistry.
Chemical language models. Chemical language models are machine
learning methods that can handle molecular sequences as inputs
and/or outputs. The most common algorithms for chemical lan-
guage modelling are RNNs and transformers:
•
RNNs (Fig. 3b)94 are neural networks that process sequence data
as Euclidean structures, usually via one-hot-encoding. RNNs
model a dynamic system in which the hidden state (ht) of the
network at any tth time point (that is, at any tth position in the
sequence) depends on both the current observation (st) and the
previous hidden state (ht−1). RNNs can process sequence inputs
of arbitrary lengths and provide outputs of arbitrary lengths.
RNNs are often used in an ‘auto-regressive’ fashion, that is, to
predict the probability distribution over the next possible ele-
ments (tokens) at the time step t + 1, given the current ht and the
preceding portions of the sequence. Several RNN architectures
have been proposed to solve the gradient vanishing or exploding
problems of ‘vanilla’ RNNs95, such as long short-term memory96
and gated recurrent units97.
•
Transformers (Fig. 3c) process sequence data as non-Euclidean
structures, by encoding sequences as either (1) a fully connected
graph or (2) a sequentially connected graph, where each token is
only connected to the previous tokens in the sequence. The for-
mer approach is often used for feature extraction in general (for
example, in a transformer-encoder), whereas the latter is used for
next-token prediction (for example, in a Transformer-decoder).
The positional information of tokens is usually encoded by posi-
tional embedding or sinusoidal positional encoding6. Trans-
formers combine graph-like processing with the so-called
attention layers. Attention layers allow transformers to focus
on (‘pay attention to’) the perceived relevant tokens for each
prediction. Transformers have been particularly successful in
sequence-to-sequence tasks, such as language translation.
Extending early studies13,98, RNNs for next-token prediction
have been routinely applied to the de novo generation of molecules
with desired biological or physicochemical properties, in combina-
tion with transfer13,99,100 or reinforcement learning101,102. In this con-
text, RNNs have shown remarkable capability to learn the SMILES
syntax13,99, and capture high-level molecular features (‘semantics’),
such as physicochemical13,99 and biological properties98,100,103. In
this context, data augmentation based on non-canonical SMILES
string enumeration104,105 or bidirectional learning106 has efficiently
improved the quality of the chemical language learned by RNNs.
Most published studies have used SMILES strings or derivative
representations. One-letter amino acid sequences have been used
for peptide design107–109. RNNs have also been applied to predict
ligand–protein interactions and the pharmacokinetic properties of
drugs110,111, protein secondary structure112 and the temporal evolu-
tion of molecular trajectories113. RNNs for molecular feature extrac-
tion114,115 outperformed both traditional molecular descriptors and
graph-convolution methods for virtual screening and property
prediction in certain benchmarks114. The Fréchet ChemNet dis-
tance116, which is based on the physicochemical and biological fea-
tures learned by an RNN model, has become the de facto reference
method to capture molecular similarity in this context.
Molecular transformers have also been applied to CASP, which
can be cast as a sequence-to-sequence translation task in which the
string representations of the reactants are mapped to those of the cor-
responding product, or vice versa. Since their initial applications117,
transformers have been used to predict multi-step syntheses118, regio-
and stereoselective reactions119, enzymatic reaction outcomes120,
reaction yields and classes121,122 and other molecular properties123,124.
Transformers have also been used for de novo molecule design
by learning to translate the target protein sequence into SMILES
strings of the corresponding ligands125. Furthermore, transform-
ers have recently been combined with E(3)- and SE(3)-equivariant
layers to learn the 3D structures of proteins from their amino
acid sequences5,126. These equivariant transformers achieved
state-of-the-art performance in protein structure prediction.
a
b
c
Graph
Residual attention blocks
Feature labelling
Sequence
Feature updates
st+1
=
h
h1
h2
hT
t
st
s1
s2
s3
sT
s3
s1
s2
s5
s4
s6
s7
sT
s1
s2
st
sT–1
s2
s3
sT
pi  ∈    dυ
T
vi  ∈    dυ
CC1(C)[C@H] (C(O)=O)N2[C@@H](CC2) S1
N12[C@@H](CC1)SC(C)(C)[C@@H]2 (C(O)=O)
O=C(O)C [C@@H]1N2[C@@H](CC2) SC1(C)(C)
Fig. 3 | Chemical language modelling. a, SMILES strings, in which atom types are represented by their element symbols and bond types and branching
are indicated by other predefined alphanumeric symbols. For each molecule, via the SMILES algorithm, a string of T symbols (tokens) is obtained
(s = \{s1, s2, …, sT\}), which encodes the molecular connectivity, herein illustrated via the colour that indicates the corresponding atomic position in the graph
(left) and string (right). A molecule can be encoded via different SMILES strings depending on the chosen starting non-hydrogen atom and the direction of
string construction along the molecular graph. Three SMILES strings obtained from different starting atoms and capturing identical molecular information
are presented. b, Recurrent neural networks at any sequence position, t, learn to predict the next token st+1 of a sequence s given the current sequence
(\{s1, s2, …, st\}) and hidden state ht. c. Transformer-based language models in which the input sequence is structured as a graph are shown. Vertices are
featurized according to their token identity (for example, via token embedding, vi ∈Rdv) and their position in the sequence (for example, via sinusoidal
positional encoding, pi ∈Rdv). During transformer learning, the vertices are updated via residual attention blocks. After passing T attention layers, an
individual feature representation sT
t  for each token is obtained.
Nature Machine Intelligence | VOL 3 | December 2021 | 1023–1032 | www.nature.com/natmachintell
1028
Review Article
NaTuRe MachIne In TeLLIGence
Other deep learning approaches rely on string-based repre-
sentations for de novo design; for example, conditional generative
adversarial networks127 and variational autoencoders128. Most of
these models, however, seem to have limited or equivalent abili-
ties to automatically learn SMILES syntax compared with RNNs.
One-dimensional CNNs129,130 and self-attention networks131 have
been used with SMILES for property prediction. Recently, deep
learning on amino acid sequences for property prediction was shown
to perform on par with approaches based on human-engineered
features132.
Conclusions and outlook
GDL in chemistry has allowed researchers to leverage the symme-
tries of unstructured molecular representations, thereby increasing
the flexibility and versatility of computational models for molecu-
lar structure generation and property prediction. Such approaches
complement chemoinformatics that are based on molecular
descriptors or other human-engineered features. For modelling
tasks that are usually characterized by the need for engineered rules
(for example, molecule construction for de novo design and reactive
site specification for CASP), GDL extends the existing methodolog-
ical repertoire. In published applications of GDL, each molecular
representation has shown characteristic strengths and weaknesses.
Molecular strings, such as SMILES, have proved particularly
suited to generative deep learning tasks, such as de novo design
and CASP. This success may be attributed to the straightforward
syntax of this chemical language, which facilitates next-token and
sequence-to-sequence prediction. For molecular property predic-
tion, SMILES strings might be limited due to their non-univocity.
Molecular graphs have proved useful for property predic-
tion, partly because of their human interpretability and the ease
of inclusion of desired edge and node features. The incorpora-
tion of 3D information (for example, with E(3)-invariant or
SE(3)/E(3)-equivariant message passing) facilitates quantum
chemistry-related modelling, whereas in drug discovery appli-
cations this approach has often failed to clearly outbalance the
increased complexity of the model. E(3)-equivariant GNNs have
also been applied to conformation-aware de novo design133,134, but
are awaiting experimental validation.
Molecular grids have long been the standard 3D representa-
tion for learning tasks on large and static molecular systems, such
as proteins. Their ability to capture information at a user-defined
resolution (voxel density) and the Euclidean structure of the input
grid render 3D CNNs efficient and applicable to proteins and other
macromolecules. However, recent advances in Transformer net-
works5,126, GNNs11 and geodesic CNNs12 have led to models that
achieved state-of-the-art performance.
Finally, molecular surfaces are at the forefront of GDL at pres-
ent. We expect many interesting applications of GDL on molecular
surfaces in the near future.
To further the application and impact of GDL in chemistry,
an evaluation of the optimal trade-off between algorithmic com-
plexity, performance and model interpretability will be required.
These aspects are crucial for reconciling the ‘two QSARs’135 and
connecting the computer science and chemistry communities. We
encourage GDL practitioners to include aspects of interpretability
in their models (for example, via explainable AI43) whenever pos-
sible, and transparently communicate with domain experts. The
feedback from domain experts will also be crucial to develop new
‘chemistry-aware’ architectures and enable concrete prospective
applications.
The potential of GDL for molecular feature extraction has not
yet been fully explored. Several studies have shown the benefits of
learned representations compared with classical molecular descrip-
tors, but in other cases GDL failed to live up to its promise in terms
of superior learned features. Deriving useful data-driven features for
downstream applications might be challenging, as it requires expert
knowledge of both the algorithms and the respective domains of
application, favouring interdisciplinary collaborations. Although
there are benchmarks for evaluating machine learning models for
property prediction136,137 and molecule generation138,139, at present
there is no such framework to enable the systematic evaluation of
the usefulness of data-driven features learned by AI. Such bench-
marks and systematic studies, including prospective applications,
Box 3 | | Structure–activity landscape visualization with GDL
This example shows how GDL can be used to interpret the struc-
ture–activity landscape learned by a trained model. Starting from
a publicly available molecular dataset containing oestrogen recep-
tor binding information 149, we trained an E(3)-equivariant graph
neural network (six hidden layers, 128 hidden neurons per layer)
and analysed the learned features and their relationship to ligand
binding to the oestrogen receptor. The figure shows an analysis of
the learned molecular features (third hidden layer, analysed via
principal component analysis; the first two principal components
are shown), and how these features relate to the density of active
and inactive molecules in the chemical space. The features en-
coded in this particular intermediate network layer captured both
experimental bioactivity and structural molecular features (for ex-
ample, atom scaffolds150).
O
HO
OH
S
HO
OH
O
O
N
OH
HO
F
O
S
O
N
OH
HO
..
High
Low
Atom scaffold
Density of actives
Nature Machine Intelligence | VOL 3 | December 2021 | 1023–1032 | www.nature.com/natmachintell
1029
Review Article
NaTuRe MachIne In TeLLIGence
are indispensable for obtaining an unvarnished assessment of deep
representation learning. Moreover, investigating the relationships
between the learned features and the physicochemical and biologi-
cal properties of the input molecules will augment the interpret-
ability and applicability of GDL to modelling structure–function
relationships such as structure–activity landscapes (Box 3), for
example.
Compared with conventional QSAR approaches in which the
assessment of the applicability domain (that is, the region of the
chemical space where model predictions are considered reliable)
has been routinely performed, contemporary GDL studies lack such
an assessment. This apparent gap might constitute one of the lim-
iting factors on the more widespread use of GDL approaches for
prospective studies. Thorough assessment of a model’s applicability
domain will help reduce the risk of unreliable predictions, for exam-
ple, for molecules with mechanisms of action, functional groups or
physicochemical properties that are different from the training data.
It is advisable to develop ‘geometry-aware’ approaches for applica-
bility domain assessment.
Another opportunity will be to leverage less-explored molecu-
lar representations for GDL. For instance, the electronic structure
of molecules has potential for tasks such as CASP, molecular prop-
erty prediction and the prediction of macromolecular interactions
(for example, protein–protein interactions). Although accurate sta-
tistical and quantum mechanical simulations are computationally
expensive, modern quantum machine learning models140,141 trained
on large quantum data collections142,143 allow quantum chemical
properties to be accessed much faster with high accuracy. This
aspect could enable quantum and electronic featurization of exten-
sive molecular datasets to be used as input molecular representa-
tions for the task of interest.
Deep learning can be applied to a multitude of biological and
chemical representations. The corresponding deep neural net-
work models have the potential to augment human creativity,
paving the way for scientific studies that were previously unfea-
sible. However, research has only explored the tip of the iceberg.
One of the most important catalysts for the integration of deep
learning in molecular sciences may be the responsibility of aca-
demic institutions and other organizations to foster interdisci-
plinary education, collaboration and communication. Picking
the ‘high-hanging fruits’ will only be possible with a deep under-
standing of both chemistry and computer science, along with
out-of-the-box thinking and collaborative creativity. In such a
setting, we expect molecular GDL to increase our understanding
of molecular systems and biological phenomena.
Received: 23 July 2021; Accepted: 26 October 2021;
Published online: 15 December 2021
References
	1.
LeCun, Y., Bengio, Y. \& Hinton, G. Deep learning. Nature 521,
436–444 (2015).
	2.
Schmidhuber, J. Deep learning in neural networks: an overview. Neur. Netw.
61, 85–117 (2015).
	3.
Gawehn, E., Hiss, J. A. \& Schneider, G. Deep learning in drug discovery.
Mol. Inform. 35, 3–14 (2016).
	4.
Gilmer, J., Schoenholz, S. S., Riley, P. F., Vinyals, O. \& Dahl, G. E. Neural
message passing for quantum chemistry. In International Conference on
Machine Learning Vol. 34, 1263–1272 (2017).
	5.
Jumper, J. et al. Highly accurate protein structure prediction with
AlphaFold. Nature 596, 583–589 (2021).
	6.
Vaswani, A. et al. Attention is all you need. In Advances in Neural
Information Processing Systems Vol. 30, 5998–6008 (2017).
	7.
Hinton, G. et al. Deep neural networks for acoustic modeling in speech
recognition: the shared views of four research groups. IEEE Signal Proces.
Mag. 29, 82–97 (2012).
	8.
Krizhevsky, A., Sutskever, I. \& Hinton, G. E. Imagenet classification with
deep convolutional neural networks. In Advances in Neural Information
Processing Systems Vol. 25, 1097–1105 (2012).
	9.
Bronstein, M. M., Bruna, J., LeCun, Y., Szlam, A. \& Vandergheynst, P.
Geometric deep learning: going beyond euclidean data. IEEE Signal Process.
Mag. 34, 18–42 (2017).
	10.
Todeschini, R. \& Consonni, V. Molecular Descriptors for Chemoinformatics
Vols I–II (John Wiley \& Sons, 2009).
	11.
Townshend, R. J. et al. Geometric deep learning of RNA structure. Science
373, 1047–1051 (2021).
	12.
Gainza, P. et al. Deciphering interaction fingerprints from protein molecular
surfaces using geometric deep learning. Nat. Methods 17, 184–192 (2020).
	13.
Segler, M. H., Kogej, T., Tyrchan, C. \& Waller, M. P. Generating focused
molecule libraries for drug discovery with recurrent neural networks. ACS
Centr. Sci. 4, 120–131 (2018).
	14.
Bronstein, M. M., Bruna, J., Cohen, T. \& Veličković, P. Geometric deep
learning: grids, groups, graphs, geodesics, and gauges. Preprint at https://
arxiv.org/abs/2104.13478 (2021).
	15.
Mumford, D., Fogarty, J. \& Kirwan, F. Geometric Invariant Theory Vol. 34
(Springer Science \& Business Media, 1994).
	16.
Cohen, T. S. \& Welling, M. Group equivariant convolutional networks. In
International Conference on Machine Learning Vol. 33, 2990–2999 (2016).
	17.
Kondor, R. \& Trivedi, S. On the generalization of equivariance and
convolution in neural networks to the action of compact groups. In
International Conference on Machine Learning Vol. 35, 2747–2755 (2018).
	18.
Moriguchi, I., Hirono, S., Liu, Q., Nakagome, I. \& Matsushita, Y. Simple
method of calculating octanol/water partition coefficient. Chem.
Pharmaceut. Bull. 40, 127–130 (1992).
	19.
Cybenko, G. Approximation by superpositions of a sigmoidal function.
Math. Control Signals Syst. 2, 303–314 (1989).
	20.
Hoffmann, R. \& Laszlo, P. Representation in chemistry. Angew. Chem. Int.
Ed. Engl. 30, 1–16 (1991).
	21.
Kipf, T. N. \& Welling, M. Semi-supervised classification with graph
convolutional networks. In International Conference on Learning
Representations Vol 5. (2017).
	22.
Duvenaud, D. et al. Convolutional networks on graphs for learning
molecular fingerprints. In International Conference on Neural Information
Processing Systems, Vol. 28, 2224–2232 (2015).
	23.
Monti, F. et al. Geometric deep learning on graphs and manifolds using
mixture model CNNs. In Proc. IEEE Conference on Computer Vision and
Pattern Recognition 5115–5124 (2017).
	24.
Battaglia, P. et al. Interaction networks for learning about objects, relations
and physics. In Advances in Neural Information Processing Systems Vol. 29,
4502–4510 (2016).
	25.
Battaglia, P. W. et al. Relational inductive biases, deep learning, and graph
networks. Preprint at https://arxiv.org/abs/1806.01261 (2018).
	26.
Zhou, J. et al. Graph neural networks: a review of methods and
applications. AI Open 1, 57–81 (2020).
	27.
Schütt, K. T., Arbabzadah, F., Chmiela, S., Müller, K. R. \& Tkatchenko, A.
Quantum-chemical insights from deep tensor neural networks. Nat.
Commun. 8, 13890 (2017).
	28.
Klicpera, J., Groß, J. \& Günnemann, S. Directional message passing for
molecular graphs. In International Conference on Learning Representations
Vol. 8 (2020).
	29.
Feinberg, E. N. et al. PotentialNet for molecular property prediction. ACS
Central Science 4, 1520–1530 (2018).
	30.
Torng, W. \& Altman, R. B. Graph convolutional neural networks for
predicting drug-target interactions. J. Chem. Inf. Model. 59, 4131–4149
(2019).
	31.
Stokes, J. M. et al. A deep learning approach to antibiotic discovery. Cell
180, 688–702 (2020).
	32.
Somnath, V. R., Bunne, C., Coley, C. W., Krause, A. \& Barzilay, R. Learning
graph models for retrosynthesis prediction. In Advances in Neural
Information Processing Systems Vol. 34 (2021).
	33.
Coley, C. W. et al. A graph-convolutional neural network model for the
prediction of chemical reactivity. Chem. Sci. 10, 370–377 (2019).
	34.
Li, J., Cai, D. \& He, X. Learning graph-level representation for drug
discovery. Preprint at https://arxiv.org/abs/1709.03741 (2017).
	35.
Liu, K. et al. Chemi-Net: a molecular graph convolutional network for
accurate drug property prediction. Int. J. Mol. Sci. 20, 3389 (2019).
	36.
Unke, O. T. \& Meuwly, M. PhysNet: a neural network for predicting
energies, forces, dipole moments, and partial charges. J. Chem. Theory
Comput. 15, 3678–3693 (2019).
	37.
Schütt, K. T., Unke, O. T. \& Gastegger, M. Equivariant message passing for
the prediction of tensorial properties and molecular spectra. Preprint at
https://arxiv.org/abs/2102.03150 (2021).
	38.
Schütt, K. T., Sauceda, H. E., Kindermans, P.-J., Tkatchenko, A. \& Müller,
K.-R. SchNet–a deep learning architecture for molecules and materials. J.
Chem. Phys. 148, 241722 (2018).
	39.
Schütt, K., Gastegger, M., Tkatchenko, A., Müller, K.-R. \& Maurer, R. J.
Unifying machine learning and quantum chemistry with a deep neural
network for molecular wavefunctions. Nat. Commun. 10, 5024 (2019).
Nature Machine Intelligence | VOL 3 | December 2021 | 1023–1032 | www.nature.com/natmachintell
1030
Review Article
NaTuRe MachIne In TeLLIGence
	40.
Bogojeski, M., Vogt-Maranto, L., Tuckerman, M. E., Müller, K.-R. \& Burke,
K. Quantum chemical accuracy from density functional approximations via
machine learning. Nat. Commun. 11, 5223 (2020).
	41.
Yang, K. et al. Analyzing learned molecular representations for property
prediction. J. Chem. Inf. Model. 59, 3370–3388 (2019).
	42.
Axelrod, S. \& Gomez-Bombarelli, R. Molecular machine learning with
conformer ensembles. Preprint at https://arxiv.org/abs/2012.08452 (2020).
	43.
Jiménez-Luna, J., Grisoni, F. \& Schneider, G. Drug discovery with
explainable artificial intelligence. Nat. Mach. Intell. 2, 573–584 (2020).
	44.
Jiménez-Luna, J., Skalic, M., Weskamp, N. \& Schneider, G. Coloring
molecules with explainable artificial intelligence for preclinical relevance
assessment. J. Chem. Inf. Model. 61, 1083–1094 (2021).
	45.
Schnake, T. et al. Xai for graphs: explaining graph neural network
predictions by identifying relevant walks. Preprint at https://arxiv.org/
abs/2006.03589 (2020).
	46.
Sun, M., Xing, J., Wang, H., Chen, B. \& Zhou, J. MoCL: data-driven
molecular fingerprint via knowledge-aware contrastive learning from
molecular graph. In Proc. 27th ACM SIGKDD Conference on Knowledge
Discovery and Data Mining 3585–3594 (Association for Computing
Machinery, 2021); https://doi.org/10.1145/3447548.3467186
	47.
Coley, C. W., Barzilay, R., Green, W. H., Jaakkola, T. S. \& Jensen, K. F.
Convolutional embedding of attributed molecular graphs for physical
property prediction. J. Chem. Inf. Model. 57, 1757–1772 (2017).
	48.
Li, Y., Vinyals, O., Dyer, C., Pascanu, R. \& Battaglia, P. Learning deep
generative models of graphs. Preprint at https://arxiv.org/abs/1803.03324
(2018).
	49.
Simonovsky, M. \& Komodakis, N. GraphVAE: towards generation of small
graphs using variational autoencoders. In International Conference on
Artificial Neural Networks Vol. 27, 412–422 (2018).
	50.
De Cao, N. \& Kipf, T. MolGAN: An implicit generative model for small
molecular graphs. Preprint at https://arxiv.org/abs/1805.11973 (2018).
	51.
Zhou, Z., Kearnes, S., Li, L., Zare, R. N. \& Riley, P. Optimization of
molecules via deep reinforcement learning. Sci. Rep. 9, 10752 (2019).
	52.
Jin, W., Barzilay, R. \& Jaakkola, T. Junction tree variational autoencoder for
molecular graph generation. In International Conference on Machine
Learning Vol. 35, 2323–2332 (2018).
	53.
You, J., Liu, B., Ying, Z., Pande, V. \& Leskovec, J. Graph convolutional
policy network for goal-directed molecular graph generation. In Advances
in Neural Information Processing Systems Vol. 31, 6410–6421 (2018).
	54.
Jin, W., Barzilay, R. \& Jaakkola, T. Multi-objective molecule generation
using interpretable substructures. In International Conference on Machine
Learning Vol. 37, 4849–4859 (2020).
	55.
Lei, T., Jin, W., Barzilay, R. \& Jaakkola, T. Deriving neural architectures
from sequence and graph kernels. In International Conference on Machine
Learning Vol. 34, 2024-2033 (2017).
	56.
Xu, K., Hu, W., Leskovec, J. \& Jegelka, S. How powerful are graph neural
networks? Preprint at https://arxiv.org/abs/1810.00826 (2018).
	57.
Chen, Z., Chen, L., Villar, S. \& Bruna, J. Can graph neural networks count
substructures? In Advances in Neural Information Processing Systems Vol. 33,
10383–10395 (2020).
	58.
Bouritsas, G., Frasca, F., Zafeiriou, S. \& Bronstein, M. M. Improving graph
neural network expressivity via subgraph isomorphism counting. Preprint
at https://arxiv.org/abs/2006.09252 (2020).
	59.
Bodnar, C. et al. Weisfeiler and Lehman go topological: message passing
simplicial networks. In International Conference on Learning
Representations: Workshop on Geometrical and Topological Representation
Learning (2021).
	60.
Thomas, N. et al. Tensor field networks: rotation-and translation-equivariant
neural networks for 3D point clouds. Preprint at https://arxiv.org/abs/1802.
08219 (2018).
	61.
Satorras, V. G., Hoogeboom, E. \& Welling, M. E(n) equivariant graph
neural networks. Preprint at https://arxiv.org/abs/2102.09844 (2021).
	62.
Anderson, B., Hy, T. S. \& Kondor, R. Cormorant: covariant molecular
neural networks. In Advances in Neural Information Processing Systems Vol.
32, 14537–14546 (2019).
	63.
Miller, B. K., Geiger, M., Smidt, T. E. \& Noé, F. Relevance of rotationally
equivariant convolutions for predicting molecular properties. Preprint at
https://arxiv.org/abs/2008.08461 (2020).
	64.
Fuchs, F., Worrall, D., Fischer, V. \& Welling, M. SE(3)-transformers: 3D
roto-translation equivariant attention networks. In Advances in Neural
Information Processing Systems Vol. 33 (2020).
	65.
Unke, O. T. et al. SpookyNet: Learning force fields with electronic degrees
of freedom and nonlocal effects. Preprint at https://arxiv.org/abs/2105.00304
(2021).
	66.
Batzner, S. et al. SE(3)-equivariant graph neural networks for data-efficient
and accurate interatomic potentials. Preprint at https://arxiv.org/abs/2101.
03164 (2021).
	67.
Unke, O. T. et al. SE(3)-equivariant prediction of molecular wavefunctions
and electronic densities. Preprint at https://arxiv.org/abs/2106.02347 (2021).
	68.
Hermann, J., Schätzle, Z. \& Noé, F. Deep-neural-network solution of the
electronic Schrödinger equation. Nat. Chem. 12, 891–897 (2020).
	69.
Pfau, D., Spencer, J. S., Matthews, A. G. \& Foulkes, W. M. C. Ab initio
solution of the many-electron Schrödinger equation with deep neural
networks. Phys. Rev. Res. 2, 033429 (2020).
	70.
Choo, K., Mezzacapo, A. \& Carleo, G. Fermionic neural-network states for
ab-initio electronic structure. Nat. Commun. 11, 2368 (2020).
	71.
Rajan, K., Zielesny, A. \& Steinbeck, C. Decimer: towards deep learning for
chemical image recognition. J. Cheminform. 12, 65 (2020).
	72.
Cramer, R. D., Patterson, D. E. \& Bunce, J. D. Comparative molecular field
analysis (comfa). 1. effect of shape on binding of steroids to carrier
proteins. J. Am. Chem. Soc. 110, 5959–5967 (1988).
	73.
Klebe, G. in 3D QSAR in Drug Design (eds. Kubinyi, H. et al.) 87–104
(Springer, 1998).
	74.
Jiménez, J., Skalic, M., Martinez-Rosell, G. \& De Fabritiis, G. KDEEP:
protein–ligand absolute binding affinity prediction via 3d-convolutional
neural networks. J. Chem. Inf. Model. 58, 287–296 (2018).
	75.
Ragoza, M., Hochuli, J., Idrobo, E., Sunseri, J. \& Koes, D. R. Protein–ligand
scoring with convolutional neural networks. J. Chem. Inf. Model. 57,
942–957 (2017).
	76.
Jiménez, J., Doerr, S., Martinez-Rosell, G., Rose, A. S. \& De Fabritiis, G.
DeepSite: protein-binding site predictor using 3d-convolutional neural
networks. Bioinformatics 33, 3036–3042 (2017).
	77.
Ahmed, E. et al. A survey on deep learning advances on different 3d data
representations. Preprint at https://arxiv.org/abs/1808.01462 (2018).
	78.
Pfaff, T., Fortunato, M., Sanchez-Gonzalez, A. \& Battaglia, P. Learning
mesh-based simulation with graph networks. In International Conference on
Learning Representations Vol. 8 (2020).
	79.
Liu, Q. et al. OctSurf: efficient hierarchical voxel-based molecular surface
representation for protein-ligand affinity prediction. J. Molec. Graph. Model.
105, 107865 (2021).
	80.
Mylonas, S. K., Axenopoulos, A. \& Daras, P. DeepSurf: a surface-based
deep learning approach for the prediction of ligand binding sites on
proteins. Preprint at https://arxi
\end{document}
